% Figure for chapter 5, "An experiment": what the shard's thread does in
% the first few milliseconds of preemption.cc, for three kinds of --work.
% The x axis is time, one unit a millisecond.
\begin{tikzpicture}[x=22mm]
  % Work running, from #2 to #3 ms, in the lane at y #1, labeled #4.
  \newcommand\work[4]{
    \node[box, minimum height=5mm, minimum width={((#3)-(#2))*22mm}, inner sep=0,
          rounded corners=1.5pt, font=\sffamily\footnotesize, anchor=west]
      at (#2, #1) {#4};
  }
  % The heartbeat running, at #2 ms, in the lane at y #1.
  \newcommand\beat[2]{
    \fill[figcool, rounded corners=0.4mm] (#2-0.025, #1-3.5mm) rectangle (#2+0.025, #1+3.5mm);
  }
  % The heartbeat's timer expiring, at #2 ms, above the lane at y #1.
  \newcommand\expires[2]{
    \draw[figcool, fill=white, line width=0.72pt]
      (#2, #1+3.4mm) -- ++(-0.035, 1.8mm) -- ++(0.07, 0) -- cycle;
  }
  % How late the heartbeat runs: from #2 to #3 ms, above the lane at y #1.
  \newcommand\late[4]{
    \draw[figcool, |-|, line width=0.6pt, >={Bar[width=1.6mm]}]
      (#2, #1+7.2mm) -- (#3, #1+7.2mm)
      node[midway, above, font=\sffamily\footnotesize, text=figcool, inner sep=0.6mm] {#4};
  }
  % A lane's name, at y #1.
  \newcommand\lane[2]{
    \node[anchor=east, font=\ttfamily\footnotesize] at (-0.1, #1) {#2};
  }

  % plain: the work never gives up the CPU.
  \lane{0}{plain}
  \beat{0}{0}
  \node[box, minimum height=5mm, minimum width={4.43*22mm}, inner sep=0, anchor=west,
        rounded corners=1.5pt, font=\sffamily\footnotesize, fill=figfill]
    (plainwork) at (0.07, 0) {work, for the whole second, never giving up the CPU};
  \expires{0}{1}
  \draw[figcool, message, rare] (1.05, 5.4mm) -- (4.55, 5.4mm)
    node[pos=0.5, above, font=\sffamily\footnotesize, text=figcool, inner sep=0.6mm]
      {the heartbeat waits until the work ends, a second later};
  \node[note, anchor=south, inner sep=0.6mm] at (1, 7mm) {timer expires};
  \node[note, anchor=north, inner sep=0.8mm] at (0, -3.5mm) {heartbeat};

  % maybe_yield: the work yields when the quota is used up, or when the
  % timer expires; the heartbeat, made ready by the poll, waits behind it.
  \def\y{-19mm}
  \lane{\y}{maybe\_yield}
  \beat{\y}{0}
  \foreach \from in {0, 0.5, 1, 1.5, 2, 2.5, 3, 3.5, 4}
    \work{\y}{\from+0.035}{\from+0.475}{work};
  \foreach \t in {1.5, 3} \beat{\y}{\t};
  \foreach \t in {1, 2.5, 4} \expires{\y}{\t};
  \late{\y}{1}{1.5}{0.5 ms late}
  \late{\y}{2.5}{3}{0.5 ms late}
  \node[note, anchor=north, inner sep=0.8mm] at (0.5, \y-2.5mm) {poll};
  \draw[fignote, line width=0.5pt] (0.5, \y-2.5mm) -- (0.5, \y-3.5mm);

  % yield: the work yields after every computation, and the heartbeat runs
  % almost as soon as its timer expires.
  \def\y{-38mm}
  \lane{\y}{yield}
  \beat{\y}{0}
  \foreach \from/\to in {0.07/0.98, 1.06/2.0, 2.08/3.02, 3.1/4.04, 4.12/4.55} {
    \work{\y}{\from}{\to}{};
    \pgfmathsetmacro\first{\from + 0.04}
    \pgfmathsetmacro\step{\from + 0.08}
    \pgfmathsetmacro\last{\to - 0.03}
    \foreach \t in {\first, \step, ..., \last}
      \draw[figink!45, line width=0.4pt] (\t, \y-2.5mm) -- (\t, \y+2.5mm);
  }
  \foreach \t in {1.02, 2.04, 3.06, 4.08} \beat{\y}{\t};
  \foreach \t in {1, 2.02, 3.04, 4.06} \expires{\y}{\t};
  \node[note, anchor=north, inner sep=0.8mm, align=center] at (0.53, \y-2.5mm)
    {work, yielding every microsecond};

  % The time axis.
  \def\y{-47mm}
  \draw[figink] (0, \y) -- (4.6, \y);
  \foreach \t in {0, 0.5, ..., 4.5}
    \draw[figink] (\t, \y) -- (\t, \y-1mm);
  \foreach \t in {0, 1, 2, 3, 4}
    \node[font=\sffamily\footnotesize, anchor=north, inner sep=0.8mm] at (\t, \y-1.5mm) {\t};
  \node[font=\sffamily\footnotesize, anchor=north west, inner sep=0.8mm] at (4.6, \y-1.5mm) {ms};
\end{tikzpicture}
